% !TEX root = ../Arbeit.tex
% Gliederung der Arbeit -- die kursiven Stichpunkte unter jeder Überschrift
% beschreiben den geplanten Inhalt und werden beim Schreiben ersetzt.
\newcommand{\geplant}[1]{\textit{Geplanter Inhalt: #1}}

\chapter{Einleitung}
\label{sec:Einleitung}

\section{Ausgangssituation in der Platinenfertigung}
\label{sec:Ausgangssituation}
Die ZOPF Energieanlagen GmbH (im Folgenden ZOPF) ist ein Hersteller und
Instandhalter von Leistungselektronik für die Industrie mit Schwerpunkt auf
erneuerbaren Energien und Windkraft. Neben der Entwicklung, der Fertigung und dem
Service von Phasenmodulen, Power Stacks und laminierten Stromschienen ist ZOPF
auch in der Planung, dem Bau und der Instandhaltung von Windkraftanlagen tätig.
Zur eigenen Fertigung gehören Leiterplattenbaugruppen, also mit elektronischen
Bauelementen bestückte Leiterplatten, die im Folgenden kurz als \emph{Platinen}
bezeichnet werden. An dem Hauptstandort Leipzig werden die Platinen
gefertigt, geprüft und sowohl in Leipzig als auch in dem Zweitstandort Wickede gelagert 
und zum Teil zu größeren Baugruppen, den \emph{Modulen}, verbaut. 
Die Lackierung der Platinen übernimmt meist ein externer Dienstleister, zum Teil
erfolgt sie auch bei ZOPF selbst. 
Jede Platine durchläuft eine feste Folge von Arbeitsschritten: 
die Bestückung mit oberflächenmontierten Bauelementen
(Surface-Mounted Device, kurz SMD), die Bestückung mit bedrahteten Bauelementen
(Through-Hole Technology, kurz THT), die Prüfung mit optischer Kontrolle, die
Lackierung, die Endkontrolle und schließlich die Einlagerung. Diese Schritte
werden in Abschnitt~\ref{sec:Fertigungsprozess} näher beschrieben. Zur
Kennzeichnung trägt jede Platine einen Aufkleber mit dem Zopf-internen Namen der Platine, der Artikelnummer 
und einer \emph{Seriennummer}, welche innerhalb des Platinentyps eindeutig ist.
Auf dem Aufkleber ist zudem ein maschinenlesbarer QR-Code aufgebracht, der diese Daten codiert.

Gesteuert wird die Fertigung über \emph{Laufkarten}. Eine Laufkarte beschreibt
einen Fertigungsauftrag über eine bestimmte Stückzahl eines Platinentyps und
legt fest, welcher Bereich von Seriennummern für diesen Auftrag vergeben wird.
Bisher wurden die Laufkarten auf Papier geführt und begleiteten die Platinen
durch die Fertigung. Nach Abschluss eines Arbeitsschritts notierten die
Mitarbeiter von Hand, für welchen Seriennummernbereich der Schritt erledigt
war, und gaben diese Angaben an die Produktionsleitung weiter. Dort wurden sie
in Excel-Listen übertragen, je Platinentyp eine Arbeitsmappe mit einem
Tabellenblatt je Arbeitsschritt, und per E-Mail verteilt. Auch der Fortschritt
einer Laufkarte wurde auf Papier fortgeschrieben. Für die Module galt ein
ähnliches Verfahren: Zu jedem Modul wurde ein Montageprotokoll ausgedruckt, auf
dem die Seriennummern der eingebauten Platinen handschriftlich eingetragen
wurden. Eine digitale Erfassung dieser Zuordnung gab es nicht.

Dieses Vorgehen weist mehrere Schwächen auf. Jede Angabe wurde mindestens
einmal von Hand geschrieben und einmal übertragen; Verschreiben, Zahlendreher
und unleserliche Handschrift führten zu fehlerhaften Einträgen. Einzelne
Meldungen wurden vergessen, Zettel gingen verloren, sodass die Dokumentation
lückenhaft blieb. Die Excel-Listen ließen sich zudem jederzeit nachträglich
ändern, ohne dass der ursprüngliche Stand oder die Änderung selbst festgehalten
wurde. Eine \emph{revisionssichere} Dokumentation, bei der jede Änderung
nachvollziehbar bleibt und frühere Stände erhalten sind, war damit nicht
gegeben. Die Uneinheitlichkeit der Altdaten zeigte sich deutlich bei ihrer
Übernahme in das neue System (Abschnitt~\ref{sec:Migration}).

Darunter litten vor allem zwei Aufgaben. Zum einen war der Fortschritt einer
Laufkarte, also die Frage, wie viele Platinen welchen Arbeitsschritt bereits
abgeschlossen haben, nur durch Zusammentragen mehrerer Listen und Zettel zu
ermitteln. Zum anderen ließ sich für eine einzelne Platine kaum beantworten,
wann sie welchen Schritt durchlaufen hat, wer sie bearbeitet hat, wo sie sich
gerade befindet und in welches Modul sie eingebaut wurde. Gerade diese
\emph{Rückverfolgbarkeit} ist für die Qualitätssicherung wesentlich: Tritt ein
Fehler auf, müssen die betroffenen Platinen und Module eingegrenzt werden
können (Abschnitt~\ref{sec:Rueckverfolgbarkeit}). Die Verteilung der Platinen
auf zwei Standorte und einen externen Dienstleister erschwerte die Übersicht
zusätzlich, weil Platinen zwischen diesen Orten wechseln und die Angaben dazu
an verschiedenen Stellen vorlagen.

Vor diesem Hintergrund entstand im Rahmen einer Werkstudententätigkeit bei ZOPF
die Aufgabe, die papier- und tabellenbasierte Dokumentation durch eine
datenbankgestützte Anwendung abzulösen. Die daraus abgeleiteten Ziele und
Fragestellungen dieser Arbeit beschreibt der folgende Abschnitt.

\section{Zielsetzung und Fragestellungen}
\label{sec:Zielsetzung}
Ziel dieser Arbeit ist die Entwicklung einer datenbankgestützten
Webanwendung, die Papier-Laufkarten, handschriftliche Meldungen und Excel-Listen
in der Platinenfertigung von ZOPF ablöst. Die Anwendung trägt den Namen
\emph{Produktionsverwaltung}, da sie die gesamte Platinenproduktion verwalten soll. Entsprechend dem Titel der Arbeit
umfasst das Ziel drei Teile:

\begin{itemize}
	\item \textbf{Erfassung:} Die Mitarbeiter erfassen abgeschlossene
	Arbeitsschritte direkt in der Fertigung, entweder für einen ganzen
	Seriennummernbereich oder platinenweise durch Scannen der Aufkleber.
	Ergebnisse von Prüfständen gelangen über Schnittstellen automatisch in die
	Anwendung. Die Anwendung prüft die Eingaben bereits bei der Erfassung, etwa
	ob der vorherige Arbeitsschritt abgeschlossen ist.
	\item \textbf{Auswertung der Progressionsinformationen:} Aus den einzelnen
	Erfassungen ermittelt die Anwendung jederzeit den aktuellen Fortschritt je
	Laufkarte, je Seriennummernbereich und je Platine und weist auf nahende
	Liefertermine hin.
	\item \textbf{Rückverfolgbarkeit:} Für jede Platine ist lückenlos
	nachvollziehbar, wann, durch wen und an welchem Standort sie welchen Schritt
	durchlaufen hat, in welcher Kiste sie lag und in welches Modul sie
	eingebaut oder aus welchem sie wieder ausgebaut wurde. Die zugehörigen
	Prüfprotokolle sind direkt abrufbar. Korrekturen überschreiben keine Daten,
	sondern werden revisionssicher protokolliert.
\end{itemize}

Daraus ergeben sich vier Fragestellungen, die diese Arbeit beantwortet:

\begin{enumerate}
	\item Wie muss die Erfassung gestaltet sein, damit sie im Fertigungsalltag
	schnell und mit möglichst wenigen Fehlern gelingt?
	\item Wie lässt sich aus den erfassten Einzelschritten der Fortschritt von
	Laufkarten automatisch und aktuell auswerten?
	\item Wie kann das Datenmodell sicherstellen, dass die Historie jeder
	Platine vollständig und revisionssicher bleibt, auch bei Korrekturen,
	Reparaturen, Standortwechseln, dem Einbau in Module und einem späteren
	Austausch? Dabei ist zu klären,
	welche Regeln die Datenbank selbst durchsetzt und welche der Anwendung
	überlassen bleiben.
	\item Inwieweit erfüllt die Anwendung im praktischen Einsatz die
	Anforderungen der Fertigung und ihrer Mitarbeiter?
\end{enumerate}

Die Anwendung wurde iterativ und im laufenden Betrieb entwickelt. Bereits
frühe Versionen kamen in der Fertigung zum Einsatz, und Rückmeldungen der
Mitarbeiter flossen fortlaufend in die Weiterentwicklung ein. Die erste und
zweite Fragestellung werden in den Kapiteln~\ref{sec:Konzeption}
und~\ref{sec:Implementierung} beantwortet, die dritte zusätzlich mit Blick auf
das Datenmodell in Abschnitt~\ref{sec:Datenmodell} und die Geschäftsregeln in
Abschnitt~\ref{sec:Geschaeftsregeln}. Die vierte Fragestellung beantwortet die
Evaluation in Kapitel~\ref{sec:Evaluation}, unter anderem mit einer Befragung
der Anwender.

\section{Abgrenzung}
\label{sec:Abgrenzung}
Die Produktionsverwaltung dokumentiert die Platinenfertigung und Modulmontage, sie plant und
steuert sie nicht. Sie ist daher kein vollständiges Produktionsleitsystem
(Manufacturing Execution System, kurz MES), das Fertigungsaufträge einplant,
Maschinen ansteuert und Kapazitäten verwaltet. Ebenso wenig übernimmt sie
Aufgaben des bei ZOPF vorhandenen ERP-Systems (Enterprise Resource Planning),
etwa die Auftragsabwicklung oder die Materialwirtschaft. Eine Anbindung an
dieses System ist nicht Gegenstand der Arbeit; Laufkarten werden in der
Produktionsverwaltung selbst angelegt.

Fertigungsmaschinen wie Bestückungsautomaten sind nicht angebunden, die
Arbeitsschritte erfassen die Mitarbeiter. Automatisch gelangen nur
Ergebnisse von Prüfständen in die Anwendung: Der Prüfstand für Platinen meldet
das Ergebnis jeder Prüfung, der Prüfstand für Module meldet getestete Module
samt der eingebauten Platinen, und die Prüfprotokolle beider liest die Anwendung
aus der jeweiligen Ablage. Erfasst wird jeweils, dass und wann ein Schritt
abgeschlossen wurde, nicht seine Prozess- oder Messwerte. Die
Rückverfolgbarkeit reicht bis zur einzelnen Platine und zum Modul; welche
Bauteile oder Bauteilchargen auf einer Platine verbaut sind, erfasst die
Anwendung nicht. Ebenso wenig Gegenstand der Arbeit sind der Verkauf der Module
und das Entfernen verkaufter Module aus der Datenbank.

Die Produktionsverwaltung wird im Betrieb fortlaufend weiterentwickelt. Diese
Arbeit beschreibt den Entwicklungsstand vom 30. Oktober 2026; spätere
Änderungen sind nicht Gegenstand der Arbeit.

Gegenstand der Arbeit ist die vom Verfasser entwickelte Webanwendung samt
Datenbank. Die Software des Platinen-Prüfstandes entstand in Zusammenarbeit mit
Mitarbeitern von ZOPF; behandelt wird nur ihre Schnittstelle zur
Produktionsverwaltung. Das Programm des Modul-Prüfstandes hat der Verfasser
zwar selbst entwickelt, es ist jedoch eine eigenständige Anwendung und wird
ebenfalls nur hinsichtlich seiner Anbindung betrachtet. Erste Entwürfe der
Benutzeroberfläche und des Datenbankschemas stammen aus dem Unternehmen; sie
sind an den betreffenden Stellen (Abschnitte~\ref{sec:Datenmodell}
und~\ref{sec:Oberflaeche}) als solche gekennzeichnet.

\section{Aufbau der Arbeit}
\label{sec:Aufbau}
Kapitel~\ref{sec:Grundlagen} legt die fachlichen und technischen Grundlagen:
die Rückverfolgbarkeit in der Elektronikfertigung, den Fertigungsprozess von
Platinen, relationale Datenbanken und die Sicherung ihrer Integrität, das
Webframework Svelte 5 und SvelteKit sowie die maschinenlesbare Kennzeichnung mit Barcodes
und QR-Codes. Darauf aufbauend analysiert Kapitel~\ref{sec:Anforderungsanalyse}
den bisherigen Nachweisprozess und leitet die Akteure sowie die funktionalen
und nichtfunktionalen Anforderungen an die Anwendung ab.

Den Entwurf der Lösung mit Systemarchitektur, Datenmodell, Zustandsmodell
einer Platine, Benutzeroberfläche und Berechtigungskonzept enthält
Kapitel~\ref{sec:Konzeption}. Wie dieser Entwurf umgesetzt wurde, beschreibt
Kapitel~\ref{sec:Implementierung} entlang der drei Teilziele Erfassung,
Auswertung und Rückverfolgbarkeit, ergänzt um die in der Datenbank verankerten
Geschäftsregeln sowie die Erweiterungen um Lagerverwaltung und Module. Die
Übernahme der Excel-Altdaten und die Inbetriebnahme im Unternehmensnetz folgen
in Kapitel~\ref{sec:Inbetriebnahme}.

Kapitel~\ref{sec:Evaluation} bewertet das Ergebnis anhand von Tests, der
Leistungsoptimierung, einer Befragung der Anwender und eines Abgleichs mit den
Anforderungen; zudem zeigt es, welche Anpassungen der Praxiseinsatz nach sich
zog. Kapitel~\ref{sec:Schluss} fasst die Ergebnisse zusammen, beantwortet die
Fragestellungen aus Abschnitt~\ref{sec:Zielsetzung}, bewertet die Lösung
kritisch und gibt einen Ausblick auf mögliche Weiterentwicklungen.


\chapter{Grundlagen}
\label{sec:Grundlagen}
Dieses Kapitel legt die Grundlagen, auf denen die Anforderungen, der Entwurf und
die Umsetzung der Produktionsverwaltung aufbauen. Es beginnt mit den fachlichen
Grundlagen, der Rückverfolgbarkeit und dem Fertigungsprozess von Platinen, und
behandelt anschließend die eingesetzten Techniken: relationale Datenbanken, das
Webframework SvelteKit und die maschinenlesbare Kennzeichnung.

\section{Rückverfolgbarkeit in der Elektronikfertigung}
\label{sec:Rueckverfolgbarkeit}
Die Norm DIN EN ISO 9000, die die Grundbegriffe des Qualitätsmanagements
festlegt, definiert Rückverfolgbarkeit als die \enquote{Möglichkeit, den
Werdegang, die Verwendung oder den Ort eines Objekts zu verfolgen}
\cite[S.~19]{DINENISO9000_2026}. Bei einem Produkt kann sich die
Rückverfolgbarkeit laut einer Anmerkung zu dieser Definition auf drei Aspekte
beziehen: die Herkunft von Werkstoffen und Teilen, den Ablauf der Verarbeitung
sowie die Verteilung und den Standort des Produkts nach der Auslieferung (vgl.
\cite[S.~19]{DINENISO9000_2026}). Im Englischen ist der Begriff
\emph{Traceability} gebräuchlich. Moe stützt sich auf die gleichlautende
Definition der Vorgängernorm ISO~8402 und hebt hervor, dass Rückverfolgbarkeit
auf aufgezeichneten Identifikationen beruht. Er unterscheidet zwei Reichweiten:
die Rückverfolgbarkeit entlang der gesamten Lieferkette vom Rohstoff bis zum
Verbraucher und die \emph{interne} Rückverfolgbarkeit innerhalb eines Glieds
dieser Kette, etwa der Produktion (vgl. \cite[S.~211]{moePerspectivesTraceabilityFood1998}).
Obwohl sein Beitrag der Lebensmittelherstellung gilt, betrachtet Moe die
Grundlagen der Rückverfolgbarkeit als unabhängig vom Produkt und von der Art
der Fertigung und verweist darauf, dass Hersteller von Fahrzeugen, Flugzeugen
und Raumfahrzeugen seit Langem weit entwickelte Systeme dafür einsetzen (vgl.
\cite[S.~211]{moePerspectivesTraceabilityFood1998}). Die Produktionsverwaltung
dient in diesem Sinne der internen Rückverfolgbarkeit: Sie bildet den Ablauf der
Verarbeitung ab und hält zusätzlich fest, an welchem Standort des Unternehmens
sich eine Platine befindet. Die Herkunft der Bauteile und der Verkauf der Module
sind nach Abschnitt~\ref{sec:Abgrenzung} nicht Gegenstand der Arbeit.

Moe bezeichnet Rückverfolgbarkeit als \enquote{an essential subsystem of
quality management} \cite[S.~211]{moePerspectivesTraceabilityFood1998}. Die
Anforderungen an ein Qualitätsmanagementsystem legt die Norm DIN EN ISO 9001
fest, während DIN EN ISO 9000 die zugehörigen Grundlagen und Begriffe enthält
(vgl. \cite[S.~9\,f.]{koubekDINISO90012018}). Die ISO 9001 erschien erstmals
1987, ursprünglich um die Anforderungen an Lieferanten zu vereinheitlichen, und
wird von produzierenden Unternehmen als Lieferbedingung verwendet (vgl.
\cite[S.~8]{koubekDINISO90012018}). Seit 2026 gilt eine neue Fassung, welche
die Ausgabe von 2015 ablöst (vgl. \cite[S.~3]{DINENISO9001_2026}). Ihr
Abschnitt~8.5.2 \enquote{Kennzeichnung und Rückverfolgbarkeit} stellt an eine
Organisation vier Anforderungen (vgl. \cite[S.~34]{DINENISO9001_2026}):

\begin{itemize}
	\item die Ergebnisse ihrer Prozesse zu kennzeichnen, wenn dies für die
	Konformität der Produkte notwendig ist,
	\item den Status dieser Ergebnisse während der gesamten Produktion in Bezug
	auf die Überwachungs- und Messanforderungen zu kennzeichnen,
	\item ihre eindeutige Kennzeichnung zu steuern, wenn Rückverfolgbarkeit
	gefordert ist,
	\item die dafür notwendige dokumentierte Information als Nachweis verfügbar
	zu halten.
\end{itemize}

\noindent Der informative
Anhang der Norm erläutert den Zusammenhang: \enquote{Kennzeichnung und
Rückverfolgbarkeit sind eng miteinander verbunden, denn die Rückverfolgbarkeit
erfordert eine eindeutige Kennzeichnung der Outputs, wenn es notwendig ist,
frühere Verarbeitungsstufen oder den Status nach der Auslieferung zu bestimmen}
\cite[S.~50]{DINENISO9001_2026}. Gefordert sein kann Rückverfolgbarkeit von der
Organisation selbst, von Kunden oder durch gesetzliche und behördliche Vorgaben
(vgl. \cite[S.~50]{DINENISO9001_2026}). Besondere Bedeutung hat sie etwa bei
Medizinprodukten, Lebensmitteln und in der Automobilindustrie (vgl.
\cite[S.~100]{koubekDINISO90012018}).

Grundlage jeder Rückverfolgbarkeit von Produkten ist nach Moe, sie eindeutig
identifizieren zu können, entweder durch eine Kennzeichnung am Produkt oder an
seiner Verpackung oder durch Aufzeichnungen. Ein Rückverfolgungssystem
beschreibt er durch zwei Elemente: die \emph{Route}, also den Weg, auf dem ein
Produkt durch Fertigung und Vertrieb identifizierbar bleibt, und den
\emph{Umfang}, also die Auswahl der erfassten Angaben. Zu verfolgen sind sowohl
die Produkte als auch die Aktivitäten, die an ihnen ausgeführt werden, etwa
Verarbeitung, Lagerung und Transport. Die kleinste Einheit, die ein solches
System unterscheidet, heißt \emph{Traceable Resource Unit} (TRU). Was als TRU
gilt, legt der Entwickler des Systems fest; werden mehrere TRUs zusammengeführt,
entsteht eine neue TRU, die eine eigene Identifikation erhalten muss (vgl.
\cite[S.~212]{moePerspectivesTraceabilityFood1998}). Yilmaz et al. fassen zwei
weitere Sichtweisen zusammen: Nach Aung und Chang muss ein wirksames System die
Informationen mit dem physischen Fluss der Produkte verbinden, etwa über
Verpackungen und Beschriftungen, und Regattieri et al. nennen als seine vier
Bestandteile die Identifikation des Produkts, die Verfolgung der Daten, den
Materialfluss und ein Werkzeug zur Rückverfolgung (vgl.
\cite[S.~13]{yilmazTransparenzUndRueckverfolgbarkeit2022}).

Übertragen auf die Platinenfertigung ergibt sich das Bild in
Abbildung~\ref{fig:Rueckverfolgbarkeit}. Jede Platine ist durch ihre
Seriennummer eine eigene TRU; die Rückverfolgbarkeit reicht damit bis zum
einzelnen Stück und nicht nur bis zum Fertigungsauftrag. Die Aktivitäten sind
die Arbeitsschritte der Fertigung. Die Aufzeichnung über einen abgeschlossenen
Arbeitsschritt an einer Platine wird in dieser Arbeit als
\emph{Produktionsnachweis} bezeichnet; sie hält fest, wann, durch wen und an
welchem Standort der Schritt erledigt wurde. Werden mehrere Platinen in ein
Modul eingebaut, entsteht im Sinne Moes eine neue Einheit mit eigener
Seriennummer, über die die eingebauten Platinen weiterhin auffindbar bleiben
müssen. Regulär werden dabei nur Platinen aus dem Lager eingebaut, die also die
Fertigung abgeschlossen haben; die Ausnahmen beschreibt
Abschnitt~\ref{sec:Module}.

\begin{figure}[htbp]
	\centering
	\begin{tikzpicture}[
		font=\footnotesize,
		objekt/.style={draw, rounded corners=2pt, fill=black!5, align=center,
			text width=27mm, minimum height=9mm},
		schritt/.style={draw, fill=white, align=center, minimum width=13mm,
			minimum height=8mm, inner sep=2pt},
		pfeil/.style={-{Stealth[length=2mm]}},
		]
		\node[objekt, dashed, fill=none, text=black!60] (bauteile) at (0,0.6)
			{Bauteile und\\Bauteilchargen};
		\node[right=4mm of bauteile, text=black!60, align=left]
			{nicht Gegenstand der Arbeit\\(Abschnitt~\ref{sec:Abgrenzung})};
		\node[objekt, double copy shadow={shadow xshift=2pt, shadow yshift=-2pt}]
			(platine) at (0,-2) {Platinen\\(je eine\\Seriennummer)};
		\draw[pfeil, dashed, black!60] (bauteile) -- (platine);

		\node[schritt, right=8mm of platine] (s1) {SMD};
		\node[schritt, right=3mm of s1] (s2) {THT};
		\node[schritt, right=3mm of s2] (s3) {Prüfung};
		\node[schritt, right=3mm of s3] (s4) {Lackierung};
		\node[schritt, right=3mm of s4] (s5) {End-\\kontrolle};
		\node[schritt, right=3mm of s5] (s6) {Lager};
		\draw (platine) -- (s1);
		\foreach \a/\b in {s1/s2, s2/s3, s3/s4, s4/s5, s5/s6}
			\draw[pfeil] (\a) -- (\b);
		\node[above=1mm of s3.north east] (titel)
			{Produktionsnachweise je Arbeitsschritt};
		\node[below=1mm of s3.south east, align=center] (merkmale)
			{jeweils mit Zeitpunkt, Mitarbeiter und Standort};

		\node[objekt, below=20mm of s6] (modul) {Modul\\(eigene\\Seriennummer)};
		\draw[pfeil] (s6) -- node[pos=0.55, left=1mm] {Einbau} (modul);

		\node[draw, dashed, rounded corners=4pt, inner sep=3mm,
			fit=(platine)(modul)(s6)(titel)(merkmale),
			label={[anchor=south west, font=\footnotesize\itshape]south west:Umfang der Produktionsverwaltung}]
			(umfang) {};

		\node[objekt, dashed, fill=none, text=black!60, anchor=north]
			(verkauf) at ($(modul |- umfang.south)+(0,-7mm)$)
			{Verkauf des\\Moduls};
		\node[left=4mm of verkauf, text=black!60, align=right]
			{nicht Gegenstand der Arbeit\\(Abschnitt~\ref{sec:Abgrenzung})};
		\draw[pfeil, dashed, black!60] (modul) -- (verkauf);
	\end{tikzpicture}
	\caption{Objekte und Aufzeichnungen der Rückverfolgbarkeit in der Platinenfertigung}
	\quelle{Eigene Darstellung in Anlehnung an \cite[S.~212]{moePerspectivesTraceabilityFood1998}}
	\label{fig:Rueckverfolgbarkeit}
\end{figure}

Damit eine solche Rückverfolgung im Bedarfsfall gelingt, müssen die
Aufzeichnungen bestimmten Anforderungen genügen. Ein Produktionsnachweis ist im
Sinne der DIN EN ISO 9000 eine \emph{Aufzeichnung}, also ein \enquote{Dokument,
das erreichte Ergebnisse angibt oder einen Nachweis ausgeführter Tätigkeiten
bereitstellt} \cite[S.~26]{DINENISO9000_2026}. Als einen Zweck solcher
Aufzeichnungen nennt die Norm ausdrücklich die Darlegung von
Rückverfolgbarkeit. Aufzeichnungen sind Teil der \emph{dokumentierten
Information}, die in jedem Format und auf jedem Medium vorliegen kann (vgl.
\cite[S.~26]{DINENISO9000_2026}), auf Papier ebenso wie in einer Datenbank.
DIN EN ISO 9001 verlangt, dass dokumentierte Information verfügbar ist, wo und
wann sie benötigt wird, dass sie vor dem Verlust ihrer Integrität geschützt
wird und dass Änderungen an ihr überwacht werden, etwa durch eine
Versionskontrolle (vgl. \cite[S.~27]{DINENISO9001_2026}). Für Nachweise gilt
darüber hinaus: \enquote{Dokumentierte Information, die als Nachweis der
Konformität verfügbar ist, muss vor unbeabsichtigten Änderungen geschützt
werden} \cite[S.~28]{DINENISO9001_2026}. Der informative Anhang nennt zusätzlich
den Schutz vor unbefugten Änderungen und vor dem Löschen (vgl.
\cite[S.~48]{DINENISO9001_2026}).

Moe betont, dass Rückverfolgbarkeit rückwärts arbeitet: Die Aufzeichnungen
werden erst im Nachhinein abgefragt, etwa wenn ein Fehler aufgetreten ist, und
müssen deshalb von vornherein für diese Abfrage ausgelegt sein. Unverzichtbar
ist dafür ein stabiles und zugängliches System zur Aufzeichnung (vgl.
\cite[S.~212]{moePerspectivesTraceabilityFood1998}). Rückverfolgungssysteme
können auch auf Papier funktionieren; mit Rechnerunterstützung lassen sich
jedoch deutlich mehr Daten verarbeiten und die Geschichte eines Produkts und
seiner Verarbeitung sehr detailliert festhalten. Als Nutzen einer internen
Rückverfolgbarkeit nennt Moe unter anderem eine bessere Prozesssteuerung,
Hinweise auf Ursachen, wenn ein Produkt nicht den Vorgaben entspricht, und
eine einfachere Informationsbeschaffung bei Audits des Qualitätsmanagements
(vgl. \cite[S.~213]{moePerspectivesTraceabilityFood1998}).

Aus diesen Grundlagen ergeben sich drei Voraussetzungen für die
Produktionsverwaltung: eine eindeutige, maschinenlesbare Kennzeichnung jeder
Platine (Abschnitt~\ref{sec:Kennzeichnung}), eine lückenlose Aufzeichnung der
Arbeitsschritte, deren Ablauf Abschnitt~\ref{sec:Fertigungsprozess} beschreibt,
und eine Datenhaltung, die diese Aufzeichnungen vor unbemerkter Änderung und
Verlust schützt (Abschnitt~\ref{sec:RelationaleDatenbanken}).

\section{Fertigungsprozess von Leiterplattenbaugruppen}
\label{sec:Fertigungsprozess}
Eine Leiterplattenbaugruppe entsteht, indem elektronische Bauteile auf einer
Leiterplatte bestückt und gelötet werden (vgl.
\cite[S.~156]{zickertLeiterplatten2023}). Zickert unterscheidet dabei zwei
Montagearten, die Abbildung~\ref{fig:Montagearten} gegenüberstellt. Bei der
\emph{Durchsteckmontage} (THT) werden die Anschlüsse bedrahteter Bauelemente
durch Bohrungen in der Leiterplatte gesteckt und auf der gegenüberliegenden Seite
verlötet. \emph{Oberflächenmontierte} Bauelemente (SMD) werden dagegen auf
derselben Seite platziert und verlötet. Ihre Anschlüsse und Lötflächen können
kleiner ausfallen, und beide Seiten der Leiterplatte stehen für die Bestückung
zur Verfügung, dadurch ist eine höhere Packungsdichte möglich (vgl.
\cite[S.~65\,f.]{zickertLeiterplatten2023}). Beide Montagearten lassen sich auf
einer Leiterplatte kombinieren (vgl. \cite[S.~92]{zickertLeiterplatten2023}).
Die Oberflächenmontage ersetzt die Durchsteckmontage zwar zunehmend (vgl.
\cite[S.~66]{zickertLeiterplatten2023}), bedrahtete Bauformen, Steckverbinder
und viele Bauteile der Leistungselektronik eignen sich jedoch nicht für das
Reflowlöten (vgl. \cite[S.~167]{zickertLeiterplatten2023}). Die Platinen von
ZOPF verbinden beide Montagearten, weshalb die Fertigung zwei
Bestückungsschritte umfasst.

\begin{figure}[htbp]
	\centering
	\begin{tikzpicture}[
		font=\footnotesize,
		beschriftung/.style={align=center},
		hinweis/.style={-{Stealth[length=1.5mm]}, black!60},
		]
		% Durchsteckmontage
		\node[font=\footnotesize\bfseries] at (2.5,2.3) {Durchsteckmontage (THT)};
		\fill[black!15] (0,0) rectangle (5,0.4);
		\fill[white] (1.15,-0.01) rectangle (1.45,0.41);
		\fill[white] (3.55,-0.01) rectangle (3.85,0.41);
		\draw (0,0) -- (1.15,0) (1.45,0) -- (3.55,0) (3.85,0) -- (5,0);
		\draw (0,0.4) -- (1.15,0.4) (1.45,0.4) -- (3.55,0.4) (3.85,0.4) -- (5,0.4);
		\draw (0,0) -- (0,0.4) (5,0) -- (5,0.4);
		\draw (1.15,0) -- (1.15,0.4) (1.45,0) -- (1.45,0.4);
		\draw (3.55,0) -- (3.55,0.4) (3.85,0) -- (3.85,0.4);
		\foreach \x in {1.3,3.7} {
			\fill[black!55] (\x-0.4,0) .. controls (\x-0.15,-0.03) and (\x-0.08,-0.2)
				.. (\x,-0.34) .. controls (\x+0.08,-0.2) and (\x+0.15,-0.03)
				.. (\x+0.4,0) -- cycle;
		}
		\draw[very thick] (1.3,-0.34) -- (1.3,1.15) -- (1.8,1.15);
		\draw[very thick] (3.7,-0.34) -- (3.7,1.15) -- (3.2,1.15);
		\draw[fill=white, rounded corners=2pt] (1.8,0.9) rectangle (3.2,1.4);
		\node[beschriftung] at (2.5,1.75) {bedrahtetes Bauelement};
		\node at (2.5,0.2) {Leiterplatte};
		\node[beschriftung] (l1) at (2.5,-0.75) {Lötstellen auf der Gegenseite};
		\draw[hinweis] (l1.north west) ++(0.3,0) -- (1.5,-0.15);
		\draw[hinweis] (l1.north east) ++(-0.3,0) -- (3.5,-0.15);

		% Oberflächenmontage
		\begin{scope}[xshift=6.6cm]
			\node[font=\footnotesize\bfseries] at (2.5,2.3) {Oberflächenmontage (SMD)};
			\draw[fill=black!15] (0,0) rectangle (5,0.4);
			\fill[black!80] (1.25,0.4) rectangle (2.0,0.46);
			\fill[black!80] (3.0,0.4) rectangle (3.75,0.46);
			\fill[black!55] (1.3,0.46) -- (1.8,0.46) -- (1.8,0.84) -- cycle;
			\fill[black!55] (3.7,0.46) -- (3.2,0.46) -- (3.2,0.84) -- cycle;
			\draw[fill=white, rounded corners=1pt] (1.8,0.46) rectangle (3.2,0.92);
			\node at (2.5,0.69) {SMD};
			\node at (2.5,0.2) {Leiterplatte};
			\node[beschriftung] (l2) at (2.5,1.75) {Lötstellen auf derselben Seite};
			\draw[hinweis] (l2.south west) ++(0.3,0) -- (1.6,0.7);
			\draw[hinweis] (l2.south east) ++(-0.3,0) -- (3.4,0.7);
		\end{scope}
	\end{tikzpicture}
	\caption{Durchsteckmontage und Oberflächenmontage im Vergleich}
	\quelle{Eigene Darstellung in Anlehnung an \cite[S.~65]{zickertLeiterplatten2023}}
	\label{fig:Montagearten}
\end{figure}

Beide Montagearten nutzen das Löten, die häufigste und wichtigste Methode, in
der Elektronik eine elektrische Verbindung herzustellen. Dabei benetzt
geschmolzenes Lot die Verbindungspartner, die wegen des niedrigeren
Schmelzpunktes des Lotes selbst nicht anschmelzen. Flussmittel baut zuvor die
Oxidschicht auf der Kupferoberfläche ab (vgl.
\cite[S.~159\,f.]{zickertLeiterplatten2023}).

Für SMD wird das \emph{Reflowlöten} bevorzugt (vgl.
\cite[S.~169]{zickertLeiterplatten2023}). Zunächst wird Lotpaste aufgetragen,
etwa im Schablonendruck. Anschließend werden die Bauelemente bestückt, in der
Serienfertigung hauptsächlich mit Bestückungsautomaten; die leicht klebrige
Paste hält sie bis zum Löten in Position (vgl.
\cite[S.~156, 170]{zickertLeiterplatten2023}). Die Baugruppe durchläuft danach
die Zonen einer Lötanlage: In der Vorheiz- und der Haltezone wird sie erwärmt
und das Flussmittel aktiviert, in der Peakzone erfolgt das Löten, und zuletzt
kühlt sie in der Abkühlzone ab (vgl. \cite[S.~171]{zickertLeiterplatten2023}).
Typische Fehler dieses Verfahrens sind das Aufrichten eines Bauelements, der
sogenannte Grabsteineffekt, und Fehlplatzierungen durch das Aufschwimmen der
Bauelemente auf dem flüssigen Lot (vgl.
\cite[S.~172\,f.]{zickertLeiterplatten2023}).

Für bedrahtete Bauelemente eignet sich das Wellenlöten; sollen nur einzelne
Lötstellen gelötet werden, bietet sich das Selektivlöten mit einer Miniwelle an
(vgl. \cite[S.~167, 183]{zickertLeiterplatten2023}). Bei ZOPF werden sie im
Schritt THT, im Gegensatz zu SMD-Bauteilen, von Hand gelötet.

Anschließend wird die Platine geprüft. Ein Mitarbeiter testet sie am Prüfstand
und kontrolliert sie optisch auf bekannte oder offensichtliche Fehler, etwa
verkehrt herum eingelötete Bauelemente oder kalte Lötstellen. Eine automatische
optische Inspektion wird nicht eingesetzt.

Danach erhält die Platine eine Schutzlackierung. Korrosion und Feuchte sind ein
häufiger Grund für den Ausfall elektronischer Baugruppen, und eine
Schutzbeschichtung bewahrt gefährdete Bereiche vor dem direkten Kontakt mit
Umgebungsluft und Feuchte (vgl.
\cite[S.~45, 47]{suppaAnforderungenSchutzlackierung2013}). Ob und wie eine
Baugruppe beschichtet wird, ergibt sich aus ihrer Funktion und ihren
Einsatzbedingungen. Infrage kommen unter anderem das Aufsprühen, die
Tauchlackierung und eine partielle Beschichtung (vgl.
\cite[S.~44]{suppaAnforderungenSchutzlackierung2013}); Flächen wie Prüfpunkte
und Kontaktflächen müssen frei von Lack bleiben (vgl.
\cite[S.~46]{suppaAnforderungenSchutzlackierung2013}). Bei ZOPF übernimmt diesen
Schritt meist ein externer Dienstleister, zum Teil wird auch im Unternehmen
selbst lackiert. Extern lackierte Platinen verlassen dafür das Unternehmen,
weshalb auch ihr Weg dorthin und zurück rückverfolgbar sein muss
(Abschnitt~\ref{sec:Lagerverwaltung}).

Den Abschluss bildet die Endkontrolle. Die Platine wird erneut optisch geprüft,
und es wird beurteilt, ob der Lack einwandfrei ist. Als Fehlerbilder einer
Schutzbeschichtung nennt die Richtlinie IPC-A-610 unter anderem den Verlust der
Haftung, Bläschen und Einschlüsse von Fremdkörpern (vgl.
\cite[S.~45]{suppaAnforderungenSchutzlackierung2013}). Die Endkontrolle dient
bei ZOPF als Vier-Augen-Prinzip und soll ausschließen, dass in den vorherigen
Schritten etwas übersehen wurde. Nach der ISO 9001 setzt die Freigabe von
Produkten eine solche Verifizierung voraus; aufzubewahren sind Aufzeichnungen
darüber, ob die Annahmekriterien erfüllt sind und wer freigegeben hat (vgl.
\cite[S.~103]{koubekDINISO90012018}). Nach der Endkontrolle wird die Platine
eingelagert.

Nicht jede Platine durchläuft diese Kette fehlerfrei. Für fehlerhafte Ergebnisse
sieht die ISO 9001 unter anderem die Korrektur, das Aussondern und das Sperren
vor; Aufzeichnungen über den Fehler, die Maßnahmen und die Entscheidung sind
aufzubewahren (vgl. \cite[S.~105]{koubekDINISO90012018}). Fehler werden etwa
durch Nacharbeit oder Reparatur behoben (vgl.
\cite[S.~106]{koubekDINISO90012018}). Bei ZOPF gibt es dafür neben der
Hauptkette die Sonderschritte Reparatur und Gesperrt. Nach einer Reparatur kann
eine Platine Arbeitsschritte erneut durchlaufen, sodass derselbe Schritt an einer
Platine mehrfach vorkommt (Abschnitt~\ref{sec:Zustandsmodell}). Damit die
Rückverfolgbarkeit erhalten bleibt, darf eine Wiederholung den früheren
Nachweis nicht überschreiben. Wie eine Datenbank solche Anforderungen
unterstützt, beschreibt der folgende Abschnitt.

\section{Relationale Datenbanken und Integritätssicherung}
\label{sec:RelationaleDatenbanken}
\geplant{Relationenmodell, Schlüssel und Constraints, Trigger, Transaktionen, Historisierung von Daten.}

\section{Webanwendungen mit SvelteKit}
\label{sec:SvelteKit}
\geplant{Client-Server-Architektur, serverseitiges Rendering, Load-Funktionen und Form-Actions, Svelte-5-Runes.}

\section{Maschinenlesbare Kennzeichnung}
\label{sec:Kennzeichnung}
\geplant{Barcode (Code 128) und QR-Code, Barcode-Scanner als Tastatureingabe, Etikettendruck.}


\chapter{Anforderungsanalyse}
\label{sec:Anforderungsanalyse}

\section{Analyse des bisherigen Nachweisprozesses}
\label{sec:IstAnalyse}
\geplant{Ist-Ablauf mit Laufkarten und Excel, Medienbrüche, Schwachstellen.}

\section{Akteure und Rollen}
\label{sec:Akteure}
\geplant{Gast, Mitarbeiter, Administrator; Anwendungsfälle je Rolle.}

\section{Funktionale Anforderungen}
\label{sec:FunktionaleAnforderungen}
\geplant{Erfassung von Produktionsschritten, Auswertung des Fortschritts, Rückverfolgung je Platine, Korrektur, Lager und Standorte, Modulmontage.}

\section{Nichtfunktionale Anforderungen}
\label{sec:NichtfunktionaleAnforderungen}
\geplant{Revisionssicherheit, Bedienbarkeit in der Fertigung, Antwortzeiten, Sicherheit, Wartbarkeit.}


\chapter{Konzeption}
\label{sec:Konzeption}

\section{Systemarchitektur}
\label{sec:Systemarchitektur}
\geplant{Komponentenübersicht: Browser, SvelteKit-Server, PostgreSQL, Fileserver für Prüfprotokolle.}

\section{Datenmodell}
\label{sec:Datenmodell}
\geplant{ER-Modell: Platinentypen, Laufkarten mit Seriennummernbereichen, Platinen, Produktionsnachweise mit Versuchen, Änderungsprotokoll, Standorte, Module mit Soll-Bestückung und Platinenzuordnung; nur anhängende Protokolltabellen für Kisten- und Modulzuordnungen, eingelesene Modulprüfungen; Seriennummernkonzept einschließlich Nummernversatz für Fremdfertiger-Platinen.}

\section{Zustandsmodell einer Platine}
\label{sec:Zustandsmodell}
\geplant{Hauptkette der Produktionsschritte, Sonderschritte Reparatur und Gesperrt, zweiter Durchlauf nach Reparatur mit direkter Einlagerung nach der Endkontrolle.}

\section{Gestaltung der Benutzeroberfläche}
\label{sec:Oberflaeche}
\geplant{UI-Entwürfe der Seiten Erfassung, Suche, Korrektur und Laufkartenverwaltung (Herkunft der Entwürfe kennzeichnen).}

\section{Anmelde- und Berechtigungskonzept}
\label{sec:Berechtigungskonzept}
\geplant{Rollen und Rechte, Nutzeranlage ohne Selbstregistrierung, Erstanmeldung per Startcode, Passwort-Hashing mit argon2.}


\chapter{Implementierung}
\label{sec:Implementierung}

\section{Entwicklungsvorgehen und Werkzeuge} 
\label{sec:Entwicklungsvorgehen}
\geplant{Iteratives Vorgehen, Versionsverwaltung mit Git, Feature-Branches und Pull Requests, Docker.}

\section{Erfassung von Produktionsschritten}
\label{sec:Erfassung}
\geplant{Erfassung je Seriennummernbereich, Prüfung der Vorgängerschritte, Scan-Modus für alte und neue Aufkleber, Auflösung mehrdeutiger Seriennummern eines Fremdfertigers über den Standort.}

\section{Geschäftsregeln in der Datenbank}
\label{sec:Geschaeftsregeln}
\geplant{Trigger für gültige Seriennummernbereiche und Schutz vor Bereichsverkleinerung, eindeutige aktuelle Nachweise.}

\section{Auswertung des Produktionsfortschritts}
\label{sec:Auswertung}
\geplant{Statusermittlung je Platine, Bereich und Laufkarte, Überwachung von Lieferterminen, Benachrichtigungen, Excel-Bericht je Laufkarte.}

\section{Rückverfolgbarkeit einzelner Platinen}
\label{sec:RueckverfolgbarkeitUmsetzung}
\geplant{Detailansicht mit vollständiger Historie, revisionssichere Korrektur mit Änderungsprotokoll, Korrekturanfrage, Anbindung der Prüfprotokolle, Standortverlauf, Kistenhistorie.}

\section{Lagerverwaltung und Standortwechsel}
\label{sec:Lagerverwaltung}
\geplant{Kisten und Paletten mit QR-Etikett, Versand und Wareneingang zwischen Standorten, Teillieferungen, Lagerbestandsübersicht.}

\section{Zuordnung von Platinen zu Modulen}
\label{sec:Module}
\geplant{Modultypen mit Soll-Bestückung, Seriennummernschema der Module, Montageerfassung per Scan mit Prüfung der Platinen, Modulsuche und Modul-Etikett, Austausch von Platinen mit Ein- und Ausbauprotokoll, automatisches Einlesen der Modul-Prüfprotokolle, Schnittstelle zum Modul-Prüfstand; Modulmontage bewusst kein Produktionsschritt.}


\chapter{Datenübernahme und Inbetriebnahme}
\label{sec:Inbetriebnahme}

\section{Migration der Excel-Bestandsdaten}
\label{sec:Migration}
\geplant{Import-Skripte, Bereinigung und Zuordnung der Altdaten, Problemfälle.}

\section{Bereitstellung im Unternehmensnetz}
\label{sec:Bereitstellung}
\geplant{Node-Server unter Ubuntu, PostgreSQL im Container, Einbindung des Fileservers, Datensicherung.}

\section{Einführung im Betrieb}
\label{sec:Einfuehrung}
\geplant{Anlage der Nutzer, Arbeitsanweisung, Schulung der Mitarbeiter.}


\chapter{Evaluation}
\label{sec:Evaluation}

\section{Test der Anwendung}
\label{sec:Test}
\geplant{Automatisierte Tests (Scan, Seriennummern, Prüfprotokollpfade, Module, SK-tronic-Zuordnung) und manuelle Testfälle.}

\section{Leistungsoptimierung}
\label{sec:Leistung}
\geplant{Analyse der Ladezeiten, Optimierung der Datenbankabfragen, Ergebnis (Faktor ca. 2,7).}

\section{Befragung der Anwender}
\label{sec:Befragung}
\geplant{Aufbau des Fragebogens, Durchführung in der Fertigung, Auswertung der Ergebnisse.}

\section{Erfüllung der Anforderungen}
\label{sec:Anforderungserfuellung}
\geplant{Abgleich der umgesetzten Funktionen mit den Anforderungen aus Kapitel \ref{sec:Anforderungsanalyse}.}

\section{Anpassungen aus dem Praxiseinsatz}
\label{sec:Anpassungen}
\geplant{Änderungen nach Rückmeldungen aus dem Betrieb: Reparaturregel, Aufkleberformate, Aufkleber des Fremdfertigers SK-tronic, Etikettengröße, Umbau zur Lagerverwaltung.}


\chapter{Zusammenfassung und Ausblick}
\label{sec:Schluss}

\section{Zusammenfassung der Ergebnisse}
\label{sec:Zusammenfassung}
\geplant{Kompakte Antwort auf die Fragestellungen aus Abschnitt \ref{sec:Zielsetzung}.}

\section{Kritische Bewertung}
\label{sec:KritischeBewertung}
\geplant{Grenzen der Lösung, u.\,a. clientseitige Zugriffskontrolle, fehlende automatisierte Oberflächentests, Kennzeichnung von Fremdfertiger-Platinen über einen Seriennummernversatz.}

\section{Ausblick}
\label{sec:Ausblick}
\geplant{Serverseitige Sitzungen, Module weiterer Kunden, weitere Standorte, Anbindung an das ERP-System, Anbindung von Fertigungsmaschinen.}
